% CERT rel=rec a=2 b=2 f=n T1=1 cerrada=n*log2(n)+n
\ej{11}{Árbol de recurrencia para $T(n)=2T(n/2)+n$, con $T(1)=1$.}

Suponemos $n=2^h$, con $h\in\N\cup\{0\}$, para que las divisiones lleguen exactamente al caso base.

\begin{enumerate}
\item \textbf{Construcción del árbol.} Cada nodo de tamaño $m>1$ tiene dos hijos de tamaño $m/2$ y costo propio $m$. Los tamaños de los primeros niveles son:
\begin{verbatim}
                 [n]
              /       \
          [n/2]       [n/2]
          /   \       /   \
      [n/4] [n/4] [n/4] [n/4]
        ...   ...   ...   ...
\end{verbatim}

\item \textbf{Costo por nivel.} La raíz está en el nivel $0$. En cada nivel se duplica el número de nodos y se divide entre $2$ el tamaño:
\[
\begin{array}{c|c|c|c|c}
\text{Nivel}&\text{Nodos}&\text{Tamaño}&\text{Costo por nodo}&\text{Costo total}\\ \hline
0&1&n&n&n\\
1&2&n/2&n/2&n\\
2&4&n/4&n/4&n\\
i<h&2^i&n/2^i&n/2^i&2^i(n/2^i)=n\\
h&2^h&1&T(1)=1&2^h=n
\end{array}
\]
El costo total se mantiene en $n$ por nivel; todos los niveles contribuyen por igual.

\item \textbf{Altura y hojas.} Las hojas tienen tamaño $1$, por lo que
\[
\frac{n}{2^h}=1\;\Rightarrow\;2^h=n\;\Rightarrow\;h=\log_2 n.
\]
Así, la altura es $\log_2 n$, hay $h+1=\log_2 n+1$ niveles y $2^h=n$ hojas.

\item \textbf{Costo total.} Sumamos los $h$ niveles internos y el costo de las hojas:
\[
T(n)=\sum_{i=0}^{h-1}2^i\frac{n}{2^i}+2^hT(1)
=\sum_{i=0}^{h-1}n+n
=hn+n
=n\log_2 n+n.
\]
Para $n=1$, la suma interna es vacía y la fórmula da $T(1)=1$, como se requiere.

Para $n\ge2$, se cumple $\log_2 n\ge1$, de modo que $n\le n\log_2 n$ y
\[
0\le n\log_2 n\le T(n)=n\log_2 n+n\le2n\log_2 n.
\]
Por tanto, podemos tomar $c_1=1$, $c_2=2$ y $n_0=2$.
\end{enumerate}

Concluimos que $\boxed{T(n)=n\log_2 n+n\in\Th{n\log_2 n}}$, para $n$ potencia de $2$.
\fin
